% 05 - IMPLEMENTATION
% ==============================================================
\sectionintro{05 · Structured Text}{Comptage, bacs et indexation}{Des traitements courts, directement observables en ligne dans le bloc fonctionnel.}

\noindent
\begin{minipage}[t]{0.355\textwidth}\vspace{0pt}
  \subsection{Front du capteur laser}
\begin{lstlisting}[style=stcode]
fbLaser(CLK := bLaser);

IF fbLaser.Q
AND (NOT aResetBac[nBac]) THEN
    aPieceBac[nBac] :=
        aPieceBac[nBac] + 1;
END_IF
\end{lstlisting}

  \small
  \texttt{R\_TRIG} évite d'incrémenter plusieurs fois la même pièce lorsque le signal reste actif pendant plusieurs cycles PLC. Le même test est exécuté dans \texttt{ATTENTE}, avant le passage à \texttt{COMPTAGE}, puis dans \texttt{COMPTAGE}. La condition sur \texttt{aResetBac[nBac]} donne la priorité à la réinitialisation si les deux événements surviennent dans le même cycle. Le trigger ne remplace pas un filtrage matériel ou un traitement antiparasite. Le comptage reste limité à ces états : aucune pièce ne doit être présentée avant la validation du zéro ni pendant l'indexation.

  \subsection{Données par bac}
  \begin{card}
    \texttt{aPieceBac : ARRAY[1..8] OF INT}\par\small
    Une valeur est conservée pour chaque bac. Seule la case correspondant à \texttt{nBac} est incrémentée.
  \end{card}

  \begin{notebox}{Noms retenus}
    \texttt{bLaser}, \texttt{bZero}, \texttt{nBac}, \texttt{nImp}, \texttt{tVerin} et \texttt{aResetBac} restent courts et lisibles pour une mise au point en atelier.
  \end{notebox}
\end{minipage}\hfill
\begin{minipage}[t]{0.605\textwidth}\vspace{0pt}
  \subsection{Séquence pneumatique réelle}
\begin{lstlisting}[style=stcode]
E_EtatCycle.VERIN_SORTIE:
    bVerinAv := TRUE;
    bVerinAr := FALSE;
    fbTempoVerin(IN := TRUE,
                 PT := tVerin);

    IF fbTempoVerin.Q THEN
        bVerinAv := FALSE;
        fbTempoVerin(IN := FALSE);
        eEtatCycle :=
            E_EtatCycle.VERIN_RENTREE;
    END_IF

E_EtatCycle.VERIN_RENTREE:
    bVerinAv := FALSE;
    bVerinAr := TRUE;
    fbTempoVerin(IN := TRUE,
                 PT := tVerin);

    IF fbTempoVerin.Q THEN
        bVerinAr := FALSE;
        fbTempoVerin(IN := FALSE);
        nImp := nImp + 1;

        IF nImp >= 4 THEN
            eEtatCycle :=
                E_EtatCycle.BAC_SUIVANT;
        ELSE
            eEtatCycle :=
                E_EtatCycle.VERIN_SORTIE;
        END_IF
    END_IF
\end{lstlisting}

  \begin{notebox}{Commandes déterministes}
    Au début de chaque cycle PLC, les deux commandes sont placées à \texttt{FALSE}. Seul l'état de mouvement actif commande ensuite \texttt{bVerinAv} ou \texttt{bVerinAr}. La fin d'une rentrée incrémente \texttt{nImp} sans activer de défaut.
  \end{notebox}
\end{minipage}

\newpage

% ==============================================================
